\section{Missing transverse momentum}
\label{sec:ch6_met}

The incoming partons at the LHC carry little transverse momentum, so the vector sum of the transverse momenta of all the particles is expected to be approximately zero.
An imbalance in the visible transverse momentum can therefore indicate undetected particles, such as neutrinos or weakly interacting BSM particles.
In this analysis, the neutrino produced in the LQ process and the tau-decay neutrino can together give rise to sizeable missing transverse momentum ($\mpt$).
$\mpt$ is reconstructed as the negative vector sum of the calibrated visible contributions, with its magnitude denoted by $\met$~\cite{ATLAS2025METPerformance}:
\begin{align}
    \mpt
    &=-\left(
    \sum_e\vec p_{\mathrm{T}}^{\,e}
    +\sum_\mu\vec p_{\mathrm{T}}^{\,\mu}
    +\sum_{\tau_{\mathrm{had}}}\vec p_{\mathrm{T}}^{\,\tau}
    +\sum_j\vec p_{\mathrm{T}}^{\,j}
    +\vec p_{\mathrm{T}}^{\,\mathrm{soft}}
    \right),
    \label{eq:ch6_met_vector}\\
    \met&=\left|\mpt\right|.
    \label{eq:ch6_met_magnitude}
\end{align}
Two types of contribution are included in the reconstruction:
\begin{itemize}
    \item \textbf{hard term}: The electron, muon, hadronic-tau and jet contributions selected for $\mpt$ reconstruction.
    The jet contribution is selected using the \texttt{Tight} $\mpt$ working point to suppress pile-up jets.
    This requires jet $\pT>20~\GeV$ for $|\eta|<2.5$ and $\pT>30~\GeV$ for $2.5\leq|\eta|<4.5$.
    \item \textbf{track soft term}: The vector sum of transverse momenta of selected PV tracks not used in the hard term.
\end{itemize}

\section{Object definitions in the analysis}
\label{sec:ch6_analysis_objects}

Baseline and signal objects are defined to provide two levels of selection in the analysis.
After overlap removal as described in Section~\ref{sec:ch6_overlap}, baseline objects are used for object counting and lepton vetoes.
Signal objects are selected from these baseline collections with tighter identification or isolation requirements, as summarised in Table~\ref{tab:ch6_analysis_objects}.

\begin{table}[htbp]
    \begin{spacing}{1}
    \centering
    \small
    \setlength{\tabcolsep}{4pt}
    \renewcommand{\arraystretch}{1.12}
    \caption[Baseline and signal object definitions]{Baseline object definitions and additional signal requirements used in the analysis. All baseline requirements are retained for signal objects.}
    \label{tab:ch6_analysis_objects}
    \begin{tabular}{@{}>{\raggedright\arraybackslash}p{0.14\textwidth}>{\raggedright\arraybackslash}p{0.54\textwidth}>{\raggedright\arraybackslash}p{0.26\textwidth}@{}}
        \toprule
        Object & Baseline & Additional signal requirement \\
        \midrule
        Electron & $\pT>10~\GeV$, $|\eta|<2.47$, excluding $1.37<|\eta|<1.52$.\newline
        \texttt{TightLH} identification; \texttt{Loose\_VarRad} isolation.\newline
        $|d_0|/\sigma_{d_0}<5$; $|\Delta z_0\sin\theta|<0.5~\mm$.\newline
        Calorimeter quality requirements.
        & \texttt{HighPtCaloOnly} isolation. \\
        \addlinespace[0.6em]
        Muon & $\pT>10~\GeV$, $|\eta|<2.5$.\newline
        \texttt{Loose} identification; \texttt{Loose\_VarRad} isolation.\newline
        $|d_0|/\sigma_{d_0}<3$; $|\Delta z_0\sin\theta|<0.5~\mm$.
        & \texttt{Medium} identification; \texttt{Tight\_VarRad} isolation. \\
        \addlinespace[0.6em]
        Hadronic tau & $\pT>20~\GeV$, $|\eta|<2.5$, excluding $1.37<|\eta|<1.52$.\newline
        One or three tau-decay tracks with total charge $\pm1$.\newline
        \texttt{Loose} RNN identification; \texttt{Loose} RNN electron rejection.
        & \texttt{Tight} RNN identification. \\
        \addlinespace[0.6em]
        Jet & PFlow anti-$k_t$ jets with $R=0.4$.\newline
        $\pT>20~\GeV$, $|\eta|<2.5$.\newline
        \texttt{LooseBad} cleaning; NNJVT \texttt{FixedEffPt} within its applicable kinematic range.
        & Same jet-quality requirements. \\
        \bottomrule
    \end{tabular}
    \par\smallskip
    \raggedright Jets passing the GN2v01 85\% working point are classified as $b$-jets.
    \end{spacing}
\end{table}
